\documentclass[]{article}

\usepackage{amsmath}
\usepackage{multirow}
\usepackage{natbib}
\usepackage{graphicx} 


\begin{document}

\subsection{Semi-analytic framework}

We model the statistics of the 21 cm forest using the \texttt{HAYASHI} semi-analytic framework \citep{Furlanetto_2002,Shao_2025,Xu_2009}. This code computes the distribution of neutral hydrogen absorbers along random lines of sight through the high-redshift universe, covering a redshift range from $z=7$ to $z=15$ \citep{Xu_2009,Furlanetto_2002,_oltinsk__2021,Shao_2025}. The model integrates contributions from both the diffuse intergalactic medium (IGM) and the dense gas residing within gravitationally bound minihalos, including both host halos and their subhalos \citep{Furlanetto_2002,_oltinsk__2021,Xu__2018}. For each line of sight, the framework calculates the 21 cm optical depth, $\tau$, which is inversely proportional to the gas spin temperature \citep{Xu_2009,šoltinský2025prospectsstatisticaldetection21cm}. This allows us to generate statistical distributions of absorption features, such as the cumulative number of absorbers above a certain optical depth, $N(>\tau)$, and the differential optical depth distribution, $dN/d\tau$ \citep{Furlanetto_2002,Shao_2025,šoltinský2025prospectsstatisticaldetection21cm}.

\subsection{Modeling dark matter-baryon scattering}

We incorporate the effects of elastic scattering between dark matter and baryons through two primary physical channels. The first and most significant channel is the suppression of small-scale structure \citep{Liu_2019,Dvorkin_2014,Slatyer_2018,tashiro2014effectsdarkmatterbaryonscattering,Rahimieh_2026,Halder_2021}. Momentum exchange between the two fluids damps primordial density fluctuations, leading to a cutoff in the matter power spectrum \citep{Dvorkin_2014,Liu_2019}. We model this effect by introducing a sharp cutoff in the halo mass function below a characteristic mass, $M_{cut}$ \citep{Liu_2019}. Halos with masses $M < M_{cut}$ are assumed not to form, directly reducing the population of low-mass minihalos that dominate the 21 cm forest signal \citep{Liu_2019,tashiro2014effectsdarkmatterbaryonscattering,Slatyer_2018,Halder_2021,Rahimieh_2026}.

The second channel is the thermal coupling between the two fluids. Heat transfer from the relatively warmer baryons to the colder dark matter can cool the IGM gas beyond the standard adiabatic expansion. \citep{Mu_oz_2017,short2022darkmatterbaryonscatteringeffects} We parameterize this effect with a cooling fraction, $f_{cool} = T_k / T_k^{ad}$, where $T_k$ is the kinetic temperature of the gas in the presence of scattering and $T_k^{ad}$ is the temperature from adiabatic cooling alone. Our analysis found the impact of this thermal channel on the 21 cm forest statistics to be negligible.

The cutoff mass $M_{cut}$ is directly related to the underlying dark matter-baryon scattering cross-section, $\sigma$ . We consider models where the cross-section has a power-law dependence on the relative velocity $v$, such that $\sigma = \sigma_0 v^n$ \citep{Mahdawi_2018,Rahimieh_2025}. Our analysis focuses on mapping constraints on $M_{cut}$ to the cross-section normalization $\sigma_0$ for two specific cases: a velocity-independent cross-section ($n=0$) and a Coulomb-like interaction ($n=-4$) \citep{Mahdawi_2018,Rahimieh_2025}.

\subsection{Statistical analysis and evaluation metrics}

Our analysis focuses on the differential optical depth distribution, $dN/d\tau$, as the primary observable. This distribution's shape is sensitive to the underlying population of absorbers \citep{filipp2026lsststronglensingsystems}. The suppression of low-mass halos due to dark matter-baryon scattering preferentially removes absorbers with low optical depth, imprinting a distinct "fingerprint" on the shape of the $dN/d\tau$ distribution .

To quantify the ability to distinguish a model with structure suppression from the standard $\Lambda$CDM scenario, we employ the Kullback-Leibler (KL) divergence \citep{Hee_2016,herrera2026informationtheoreticastrophysicaluncertaintieseffective}. The KL divergence, $D_{KL}(\text{CDM} \parallel M_{cut})$, measures the information gain when updating from a prior CDM distribution to a posterior distribution that includes a cutoff mass $M_{cut}$ \citep{Hee_2016,herrera2026informationtheoreticastrophysicaluncertaintieseffective}. A larger value of $D_{KL}$ indicates that the two distributions are more easily distinguishable, signifying greater intrinsic sensitivity to the effects of dark matter scattering \citep{herrera2026informationtheoreticastrophysicaluncertaintieseffective}.

\subsection{Fisher forecast methodology}

To project the observational capabilities of future radio observatories such as the Square Kilometre Array (SKA), we perform a Fisher matrix forecast \citep{Ewall_Wice_2016,Berti_2023,Balu_2023,Mason_2023}. The Fisher matrix, $F_{ij}$, quantifies the amount of information that an observable contains about a set of model parameters \citep{Mason_2023,Balu_2023,Ewall_Wice_2016,Alves_2017,sarkar2017fishermatrixbasedpredictions}. We construct the matrix for the parameter set $\{\theta_i\} = \{M_{cut}, f_{cool}, \bar{x}_{HI}\}$, where $\bar{x}_{HI}$ is the mean neutral hydrogen fraction \citep{Balu_2023,Mason_2023,Ewall_Wice_2016,zhang2025constrainingneutralhydrogenfraction}. The elements of the matrix are calculated as:
\begin{equation}
F_{ij} = \sum_{k} \frac{1}{N_k} \frac{\partial N_k}{\partial \theta_i} \frac{\partial N_k}{\partial \theta_j}
\end{equation}
where the sum is over bins of optical depth $\tau$, and $N_k$ is the number of absorbers in the $k$-th bin.\citep{Mallik_2024,2021}

A crucial element of our forecast is a realistic model for the number of available background radio sources, which we assume evolves with redshift \citep{Haiman_2004,massardi2010modelcosmologicalevolutionlow,Smol_i__2017,gao2024empiricalmodelextragalacticradio}. The number of lines of sight, $N_{los}$, is modeled as a power law:
\begin{equation}
N_{los}(z) = 10 \left( \frac{1+z}{8} \right)^{-2.5}
\end{equation}
This accounts for the increasing scarcity of bright radio-loud quasars at higher redshifts. The statistical uncertainty on the number of absorbers in each bin is assumed to be Poissonian, scaling with $1/\sqrt{N_{los}}$. By inverting the Fisher matrix, we obtain the covariance matrix, from which we extract the marginalized 1-$\sigma$ uncertainty on our parameter of interest, $\sigma(M_{cut})$ \citep{Porter_2015}. This forecast allows us to identify the optimal observational redshift window that balances intrinsic physical sensitivity with statistical constraining power \citep{wolfson2023forecastingconstraintshighzigm}.

\end{document}
                